% id: 64-31
% book: 쎈B 미적분1 · 05 도함수의 활용 (2)
% section: 유형 07 함수의 극대·극소의 활용
% figure: none
% answer: 
% answer_source: 
% variant_level: 0 (원본 전사)
% 이 파일은 build.mjs 가 items.json 에서 만든다. 고칠 때는 items.json 을 고친다.
\smash{\raisebox{4.5mm}{\dmkichul{쎈B 미적분1 64쪽 31번 · 대표 문제}}}
함수 $f(x)=\dfrac{1}{9}x^3-3x+3$에 대하여 $y=f(x)$의 그래프에서 극대가 되는 점을 $\pt{A}$, 극소가 되는 점을 $\pt{B}$라 할 때, 선분~$\pt{AB}$의 중점의 좌표는?
\begin{choices32}{$(-3{,}\mkern3mu\mathopen{}0)$}{$(0{,}\mkern3mu\mathopen{}-3)$}{$(0{,}\mkern3mu\mathopen{}3)$}{$(1{,}\mkern3mu\mathopen{}3)$}{$(3{,}\mkern3mu\mathopen{}0)$}\end{choices32}
